\documentclass[12pt,twoside]{report}

%%%%%%%%%%%%%%%%%%%%%%%%%%%%%%%%%%%%%%%%%%%%%%%%%%%%%%%%%%%%%%%%%%%%%%%%%%%%%
% Title page metadata — fill in your details
%%%%%%%%%%%%%%%%%%%%%%%%%%%%%%%%%%%%%%%%%%%%%%%%%%%%%%%%%%%%%%%%%%%%%%%%%%%%%
\newcommand{\reporttitle}{Information Structure in Spiking Neural Networks\\[0.3em]
                          with Temporal Heterogeneity}
\newcommand{\reportauthor}{Priyam Rai}
\newcommand{\supervisor}{Marcus Ghosh \& Hardik Rajpal}
\newcommand{\degreetype}{MSc Artificial Intelligence Applications and Innovation}

%%%%%%%%%%%%%%%%%%%%%%%%%%%%%%%%%%%%%%%%%%%%%%%%%%%%%%%%%%%%%%%%%%%%%%%%%%%%%
% Pass numbered citation style to natbib BEFORE includes.tex loads it
%%%%%%%%%%%%%%%%%%%%%%%%%%%%%%%%%%%%%%%%%%%%%%%%%%%%%%%%%%%%%%%%%%%%%%%%%%%%%
\PassOptionsToPackage{numbers,sort&compress}{natbib}

% Load template packages and macros
\input{includes}
\input{notation}

% Allow graphicx to find figures in the Report Figures folder (note: space in path)
% grffile is included for robust handling of spaces in directory names.
\usepackage{grffile}
\graphicspath{{../Report Figures/}{./figures/}}

\date{June 2026}

\begin{document}

% Suppress blank verso pages inserted by \cleardoublepage so they don't count
% toward the 20-page report limit (digital submission only).
\renewcommand{\cleardoublepage}{\clearpage}

% Title page
\input{titlepage}

\pagenumbering{roman}
\clearpage{\pagestyle{empty}\cleardoublepage}
\setcounter{page}{1}
\pagestyle{fancy}

%%%%%%%%%%%%%%%%%%%%%%%%%%%%%%%%%%%%
\clearpage{\pagestyle{empty}\cleardoublepage}
\fancyhead[RE,LO]{\sffamily {Table of Contents}}
\setcounter{tocdepth}{1}
\tableofcontents

\clearpage{\pagestyle{empty}\cleardoublepage}
\pagenumbering{arabic}
\setcounter{page}{1}
\fancyhead[LE,RO]{\slshape \leftmark}
\fancyhead[LO,RE]{}


%%%%%%%%%%%%%%%%%%%%%%%%%%%%%%%%%%%%
\chapter{Introduction}
%%%%%%%%%%%%%%%%%%%%%%%%%%%%%%%%%%%%

Neural networks are not designed analogues to a biological understanding of the mechanism by which a brain processes information. 
Where a standard deep learning model applies a uniform transformation at each layer, biological 
neural circuits are built from neurons that each operate at their varying tempos. Some respond 
sharply to brief bursts of input and fall silent within milliseconds; others accumulate 
evidence slowly, integrating over hundreds of milliseconds before committing to a response. 
This temporal diversity appears to be a feature rather than a flaw: it allows a population 
of neurons to simultaneously represent fast transients and slow trends, and to support 
collective coding behaviour that no individual cell could achieve alone~\cite{pereznieves2021}.

This project takes the finding of Perez-Nieves et al.~\cite{pereznieves2021} as its 
starting point. That study demonstrated clearly that SNNs with per-neuron timescale 
diversity outperform homogeneous counterparts on spoken digit classification. What it 
did not characterise is the mechanism: \emph{why}, from an information-theoretic 
perspective, does timescale diversity lead to better representations, and does it do 
so by enabling qualitatively different kinds of collective coding? Two questions follow 
directly from that gap.

The first is \emph{architectural}: does giving the neurons of a spiking neural network 
(SNN) different timescales (what we call \emph{temporal heterogeneity}) change what kinds 
of information the population can collectively encode? The second is 
\emph{methodological}: given a trained network, how do we measure whether its collective 
coding properties reflect genuine higher-order structure, rather than simply the pairwise 
correlations we would expect from the network's synchrony?

This project addresses both questions by training two distinct classes of recurrent SNNs on the 
Spiking Heidelberg Digits (SHD) dataset~\cite{cramer2022}, an auditory speech benchmark 
encoded as spike trains across 700 input channels. Two connected analytical tools are employed in 
our analysis of the information dynamics of these networks. 
The first is \emph{M-information}~\cite{liardiscalable}, an 
information-theoretic measure that captures the component of population coding that is 
genuinely synergistic: information that a group of neurons encodes together but that 
cannot be reproduced by any pair of them in isolation. M-information is assessed against a TDMI-matched
null distribution~\cite{liarinull} to control for pairwise correlation effects, yielding
a z-score that is comparable across architectures with very different correlation
structures. The second tool is principal component analysis (PCA) and manifold representations of the network's hidden
state dynamics, which reveals the geometric dimensionality of the representations
produced, a measure of how much of the representational space each network actually
uses.

By combining these lenses, the study builds a picture of how temporal heterogeneity shapes 
collective coding. The central hypothesis is that per-neuron timescale diversity will 
produce higher-dimensional representations. Our current focus will be to zero in on the 2-class binary SHD task, 
where the maximum mutual information is exactly 1 bit, allowing us to assess how efficiently each 
architecture converts input temporal structure into collective representations. 


%%%%%%%%%%%%%%%%%%%%%%%%%%%%%%%%%%%%
\chapter{Background and Literature Review}
\label{chap:bg}
%%%%%%%%%%%%%%%%%%%%%%%%%%%%%%%%%%%%

\section{Spiking Neural Networks as Temporal Processors}

Standard artificial neural networks operate on static vectors: each layer transforms 
one snapshot of activations into the next, with no intrinsic notion of time. Spiking 
neural networks take a different approach. Drawing on the mechanics of biological 
neurons, each unit in an SNN accumulates incoming current over time and fires a brief 
spike when its membrane voltage crosses a threshold, after which it resets and resumes 
integration~\cite{gerstner2002,maass1997}. Information is carried not by continuous 
activation values but by the timing and rate of these discrete events.

The most widely-used model neuron is the \emph{leaky integrate-and-fire} (LIF) unit. 
Its dynamics are governed by two coupled equations: one for the synaptic current, 
capturing how incoming spikes drive postsynaptic potentials that then decay 
exponentially, and one for the membrane potential, describing how the neuron integrates 
this current against a passive leak back towards rest. The two key parameters are the 
synaptic time constant $\tau_\text{syn}$ and the membrane time constant $\tau_\text{mem}$. 
Together, these set the neuron's effective temporal bandwidth. A small $\tau_\text{mem}$ 
means the neuron tracks rapid fluctuations in its input but forgets them quickly; a 
large $\tau_\text{mem}$ allows integration over a longer recent history at the cost of 
sensitivity to fast transients. This tuning of timescales is not merely a modelling 
parameter: it is the mechanism through which biological neurons specialise for different 
temporal features of their inputs.

Training SNNs requires differentiating through the spike-generation threshold (a
non-differentiable step function); the standard solution is surrogate gradient
descent~\cite{zenke2021}, which substitutes a smooth approximation during the backward
pass while retaining exact threshold-crossing in the forward pass.

The dataset used is the Spiking Heidelberg Digits (SHD)~\cite{cramer2022}: recordings
of spoken digits 0--9 in German and English (20 classes), encoded into spike trains
across 700 input channels via a cochlear model. SHD is a natural SNN benchmark because
the relevant information lies in spike timing patterns rather than time-averaged rates.
Our focus on SHD is also heavily inspired by the findings of Perez-Nieves et al.~\cite{pereznieves2021}, 
who demonstrated a clear performance advantage for heterogeneous SNNs on this task, 
making it an ideal testbed for our information-theoretic analysis.

\section{Temporal Heterogeneity and Its Computational Role}

A convenient simplifying assumption in many neural network models is that all neurons 
in a layer share the same parameters, including their time constants. This homogeneity 
assumption is common in standard deep learning architectures, where there are no 
timescales to speak of, and it has carried over into many SNN models by default. It is, 
however, biologically unrealistic. Recordings from cortical circuits consistently reveal 
substantial variation in intrinsic properties between neurons of the same broad type, 
including the timescales over which individual cells integrate input~\cite{gerstner2002}.

The computational consequences of this variation were investigated directly by 
Perez-Nieves et al.~\cite{pereznieves2021}, who trained recurrent SNNs on spoken 
digit classification tasks with and without neuron-level timescale diversity. Networks 
with heterogeneous time constants, where $\tau_\text{mem}$ and $\tau_\text{syn}$
varied from neuron to neuron, substantially outperformed networks using a single 
shared timescale. The intuitive explanation is that a homogeneous population has 
effectively one temporal integration channel: the same filter is applied to the input 
signal at every unit. Diverse timescales create multiple parallel channels, each 
selective for a different characteristic rate of input change. Auditory signals carry 
information simultaneously at many timescales, from the rapid onset transients that 
distinguish consonants to the slower envelope features that carry prosody and word 
identity, and a population covering a range of timescales is naturally better equipped 
to capture this full temporal richness.

In Perez-Nieves et al., heterogeneity is implemented by letting the per-neuron 
decay parameters $\alpha = e^{-\Delta t / \tau_\text{syn}}$ and 
$\beta = e^{-\Delta t / \tau_\text{mem}}$ be trained end-to-end by gradient descent, 
starting from a Gamma-distributed initialisation. This gives the network maximum 
freedom to adapt the timescale distribution to the task. The benefit of heterogeneity 
observed in that study therefore reflects a combination of two effects: the advantage 
of distributional diversity per se, and the advantage of having that distribution 
specifically tuned to the task at hand.

This project separates these two effects by introducing \emph{fitted heterogeneity}: 
the per-neuron timescales are sampled at initialisation from distributions fit to the 
empirically observed values of a pre-trained reference network, and then \emph{frozen} 
before training begins. Only the synaptic weights are subsequently updated. This creates 
a controlled setting where the diversity of timescales is fixed and consistent across 
runs, isolating the contribution of distributional form and breadth from task-specific 
tuning. Two distributional forms are tested: a log-normal fit to the empirical membrane 
timescale distribution (FittedHet-LogNorm), and a log-uniform distribution spanning the 
same empirical range (FittedHet-LogUniform). A baseline condition with a single shared 
trainable timescale (Local-Hom) provides the homogeneous reference point.


\section{Higher-Order Information in Neural Populations}
\label{sec:minfo}

\subsection{M-Information and W-Information}

The M-information and W-information framework, introduced by Liardi et al.~\cite{liardiscalable}, 
provides exactly this. For a set of neural sources $\vec{S}$ and a target variable $T$, 
the total mutual information $I(\vec{S}; T)$ is decomposed as:

\begin{equation}
    I(\vec{S}; T) \;=\; W(\vec{S}; T) \;+\; M(\vec{S}; T)
    \label{eq:mw}
\end{equation}

$W(\vec{S}; T)$ is the \emph{W-information}: the maximum amount of mutual information 
between $\vec{S}$ and $T$ that could be produced by any distribution $Q$ sharing the 
same pairwise marginals as the observed joint distribution $P$. In other words, $W$ 
captures everything pairwise relationships can explain. $M(\vec{S}; T)$ is the residual: 
information that remains once all pairwise explanations have been exhausted, and which 
therefore reflects genuine higher-order collective structure.

Computing $W$ requires minimising the mutual information over all pairwise-consistent 
distributions, which is generally intractable. Liardi et al.\ show that when $\vec{S}$ 
is modelled as jointly Gaussian (after applying a Gaussian copula transformation to
the raw data), the problem becomes a convex optimisation over the Cholesky 
decomposition of the covariance matrix, making computation tractable for populations 
of up to around 32 neurons~\cite{liardiscalable}. This is the maximum subset size 
used in this study.

The Gaussian copula normalisation maps each neuron's marginal distribution to a 
standard Gaussian while preserving the rank-order relationships between neurons. This 
step is validated for non-Gaussian data by Liardi et al.\ using the Wilson-Cowan 
neural population model: M-information correctly identifies the critical oscillatory 
regime even when the raw activity distributions are markedly non-Gaussian, justifying 
its use on both continuous membrane potentials and discrete spike trains in the present 
project~\cite{liardiscalable}. 

\subsection{Prior Work: Information-Theoretic Analysis of Spiking Networks}

Ah-Weng and Rajpal~\cite{ahweng} applied O-information and S-information~\cite{rosas2019}
to SNN spike trains, showing that information-theoretic signatures detect dynamical
regime transitions. The present project extends this by applying M-information to
\emph{trained} classification networks and adding null z-scoring for robust
cross-architecture comparison.


\section{Null Models for Significance and Comparability}
\label{sec:null}

Raw M-information values are difficult to interpret alone and compare across 
networks. A network whose neurons fire in near-synchrony (because, for example, all
neurons share the same timescale and track the input in lockstep) will naturally 
produce elevated M-information values, not because it has richer collective coding but 
because tight pairwise correlations inflate the baseline against which higher-order 
structure is measured. Comparing raw M across such a network and a more heterogeneous 
one would risk attributing to architecture a difference that is really driven by 
pairwise synchrony.

The NuMIT framework of Liardi et al.~\cite{liarinull} addresses this through systematic 
null model comparison. The core idea is to ask: given the pairwise statistics of a 
network, how much M-information would we expect if there were no higher-order structure 
at all? A null model is constructed to have exactly the observed pairwise correlations 
while being forced to have $M = 0$ by construction. The observed M is then assessed 
against a distribution of such null samples, asking how far the real network departs 
from this pairwise-only baseline.

Concretely, because \texttt{wimfo}'s discrete interface operates on 2-neuron delayed
transitions $(x_1, x_2, y_1, y_2)$, the null is constructed at the pair level.
For each neuron pair, the 2-neuron \emph{temporal delayed mutual information}
(TDMI) is measured from the binary spike sequences at a fixed 15\,ms lag and used
as the pairwise matching constraint. Null samples are drawn from a
\emph{TDMI-matched discrete family} following the NuMIT recipe~\cite{liarinull}:
source-state marginals are sampled from a Dirichlet prior, a logic gate from a
fixed gate family (XOR and OR variants) defines the deterministic transition
structure, and bounded flip noise is calibrated to reproduce the observed TDMI
exactly. Twenty null draws are generated per pair, W and M are evaluated on each
null PMF using \texttt{wimfo}'s discrete path, and the observed pair-level $M$ is
z-scored against this null $M$ distribution. For subset sizes $k > 2$, pairs
within the subset are sampled and their z-scores aggregated to produce a
subset-level summary. The z-score is:

\begin{equation}
    z \;=\; \frac{M_\text{observed} - \bar{M}_\text{null}}{\sigma_\text{null}}
    \label{eq:zscore}
\end{equation}

A large positive z-score means the network's collective coding structure substantially 
exceeds what its own pairwise correlations would predict. The z-score, rather than raw 
M, is the primary architecture comparison metric throughout this study. 


\section{Representational Geometry in Recurrent Networks}
\label{sec:geometry}

\subsection{Low-Dimensional Dynamics in High-Dimensional State Spaces}

A recurrent network with $n$ neurons formally occupies an $n$-dimensional state space, 
but trained recurrent networks rarely use all of it. Sussillo and Barak~\cite{sussillo2013} 
established a foundational result in this area: the hidden state trajectories of trained 
recurrent neural networks tend to lie near a low-dimensional manifold, and projecting 
these trajectories onto the first few principal components of the hidden activations 
reveals interpretable computational structure (fixed points, slow manifolds,
rotational dynamics) that explains how the task is being solved. This is not a 
visualisation artifact; it reflects the fact that optimisation during training pushes 
activity into a compact, task-relevant subspace, discarding the many degrees of freedom 
that are not needed.

\subsection{Variance Explained as a Quantitative Summary}

Wakhloo et al.~\cite{wakhloo2026} provide a theoretical underpinning for PCA-based
geometry analysis. They define the \emph{participation ratio} (PR) of the neural
covariance eigenspectrum as:

\begin{equation}
    \text{PR} \;=\; \frac{\left(\displaystyle\sum_i \lambda_i\right)^2}{\displaystyle\sum_i \lambda_i^2}
    \label{eq:pr}
\end{equation}

The PR is high when variance is spread evenly across many components and low when one
or two directions dominate~\cite{wakhloo2026}.

\subsection{Comparing Representational Geometries Across Networks}

Cross-architecture comparisons will use RSA~\cite{kriegeskorte2008} and
CKA~\cite{kornblith2019} to measure pairwise representational similarity,
UMAP~\cite{mcinnes2018} for class manifold visualisation, and the manifold
capacity framework of Chung et al.~\cite{chung2018} for a task-informed
geometric summary comparable to the M-information analysis.


%%%%%%%%%%%%%%%%%%%%%%%%%%%%%%%%%%%%
\chapter{Technical Achievements}
\label{chap:achiev}
%%%%%%%%%%%%%%%%%%%%%%%%%%%%%%%%%%%%

\section{Network Architectures and Design}
\label{sec:architectures}

So far, nine recurrent SNNs are trained and analysed. All share the same topology: a 
700-channel input layer feeds a single recurrent hidden layer of 32 LIF neurons, 
connected to a fully-connected readout layer. The networks differ only in how 
the hidden-layer time constants are initialised and whether they are updated 
during training:

\begin{itemize}

    \item \textbf{Local-Hom}: A single shared time constant is maintained for all 
    hidden neurons. It is trained by gradient descent alongside the synaptic 
    weights. All neurons are temporally identical, but the population converges 
    to whatever average timescale the task rewards.

    \item \textbf{FittedHet-LogNorm}: Per-neuron $\tau_\text{mem}$ values are 
    sampled at initialisation from a log-normal distribution fitted to the 
    empirically observed membrane timescale distribution of a pre-trained reference 
    network (15-epoch run). Synaptic time constants are drawn from a Gamma 
    distribution fitted to the same source. Both are frozen via 
    \texttt{requires\_grad\_(False)} before training begins; only the synaptic 
    weights are subsequently updated.

    \item \textbf{FittedHet-LogUniform}: Same procedure as FittedHet-LogNorm, 
    but $\tau_\text{mem}$ is sampled from a log-uniform distribution spanning 
    the observed empirical $[\tau_\text{min}, \tau_\text{max}]$ range. This 
    provides flatter, broader coverage of the timing spectrum compared to the 
    log-normal variant.

    \item \textbf{Repo-Learned-Het} (\emph{reference}): Per-neuron timescales 
    are initialised from a Gamma distribution and then updated by gradient 
    descent throughout training, following Perez-Nieves et al.~\cite{pereznieves2021}. 
    This network is trained only on the all-class task and serves as a reference 
    point for the fully task-adaptive heterogeneity condition. (Not directly comparable)

\end{itemize}

A key design decision is that all three local network families are trained under 
identical hyperparameters (batch size 256, learning rate $4 \times 10^{-3}$, 
25 epochs, gradient clipping at norm 1). This ensures that any observed 
differences in performance or information structure are attributable to the 
timescale design, not to differences in optimisation settings.

A 32-neuron hidden layer was adopted because the reference implementation of
Perez-Nieves et al.~\cite{pereznieves2021} uses 128 neurons, which renders
full-population M-information analysis intractable (the \texttt{wimfo} optimisation
is bounded at $k = 32$~\cite{liardiscalable}); subsampling a larger network would
introduce subset-selection bias. The same SHD dataset, cochlear encoding, LIF model,
and surrogate gradient procedure~\cite{zenke2021, cramer2022} are retained for
comparability. The Repo-Learned-Het condition replicates the Perez-Nieves training
procedure at this smaller scale, confirming the heterogeneity advantage transfers
before any fixed-distribution analysis is performed; absolute accuracy is reduced
relative to the 128-neuron network, so comparisons with published baselines are
qualitative.

\section{Task Conditions}

Three tasks are defined from SHD~\cite{cramer2022}, forming a deliberate complexity 
gradient. The all-class condition directly replicates the task studied by 
Perez-Nieves et al.~\cite{pereznieves2021} and provides the highest-complexity 
baseline. The two parity tasks extend the analysis to ask whether the heterogeneity 
advantage and the information-theoretic signatures observed at full twenty-class 
complexity are preserved when the classification problem is substantially simplified, 
or whether they are specific to the high-complexity regime.

\begin{itemize}
    \item \textbf{All-class (20-way)}: All 20 SHD classes. Requires 
    discrimination across the full spectro-temporal diversity of the corpus.

    \item \textbf{2-class parity}: Two classes from SHD (one English digit, 
    one German digit from the same numeral). Binary discrimination between two 
    temporally similar spoken words.

    \item \textbf{4-class parity-language}: Four classes, comprising two English
    digits and two German equivalents. Requires simultaneous discrimination of 
    digit identity and language, a cross-language generalisation challenge at 
    moderate complexity.
\end{itemize}

We begin with a broad sweep across varying task complexities to identify whether the expected patterns of higher-order information
are consistent across conditions. For more extensive analysis of heterogeneity we will then focus on one
task and perform detailed analyses of the information structure.

\section{Information Analysis Pipeline}

M-information and W-information are computed using the \texttt{wimfo} library 
(\texttt{W\_M\_calculator} class), which implements the Gaussian copula-based 
convex optimisation described in Liardi et al.~\cite{liardiscalable}. For each 
of the nine networks and each signal representation (membrane potential and 
spike trains), M-information is computed for randomly sampled subsets of 
$k \in \{2, 4, 8, 16, 32\}$ neurons. Sweeping over subset sizes allows 
characterisation of how collective information scales with the number of neurons 
included, revealing whether collective codes emerge gradually or require a 
critical population size to become apparent. A temporal lag of one time step 
is used throughout, capturing the one-step predictive structure of the 
hidden-layer dynamics.

The null z-scoring procedure is adopted as the primary comparison framework, as raw
M values are confounded by pairwise synchrony differences across architectures.
For each neuron pair, 20 TDMI-matched discrete null draws are generated following
the NuMIT recipe~\cite{liarinull} (Dirichlet source marginals, logic gate family,
Brentq-calibrated flip noise), and the observed pair-level M is z-scored against
this null distribution using Equation~\eqref{eq:zscore}. Results are cached to JSON.

For the PCA analysis, hidden states are mean-centred and decomposed via SVD
on the $32 \times T$ neuron-by-time matrix.

\section{Preliminary Results}

\subsection{Task Performance}

All nine networks converge to stable accuracy within the 25-epoch training 
budget, as confirmed by convergence diagnostics applied to the training curves. 
Figure~\ref{fig:accuracy} summarises final and best-epoch test accuracy across 
all networks and tasks.

\begin{figure}[htb]
    \centering
    \includegraphics[width=\textwidth]{fig01_accuracy_bar}
    \caption{Final and best-epoch test accuracy for all nine networks across 
    the three task conditions. FittedHet-LogNorm achieves the highest peak 
    accuracy in each task, with the advantage over Local-Hom being most 
    pronounced on the 4-class cross-language task.}
    \label{fig:accuracy}
\end{figure}

FittedHet-LogNorm achieves the highest accuracy on all tasks: 71.3\% (all-class),
64.4\% (4-class), and 88.2\% (2-class), compared to 65.2\%, 54.2\%, and 82.9\% for
Local-Hom. The heterogeneity advantage is largest on the 4-class task, where
FittedHet-LogNorm leads by ten percentage points above a chance level of 25\%.

\subsection{Collective Information: Raw M and Null Z-Scores}

Raw M-information values have been computed for all nine networks across both 
membrane potential and binary spike representations. However, the null 
z-scoring component of this analysis is currently undergoing a methodological 
correction and has not been finalised at the time of this report. The corrected 
pipeline is in active development and will form a central component of the 
remaining project work; see Section~\ref{sec:zscore_plan} for a full description 
of what is being corrected and why.

The qualitative hypothesis motivating the z-scoring framework remains intact: 
raw M-information values alone are insufficient for cross-architecture comparison 
because homogeneous networks tend to exhibit elevated pairwise synchrony, which 
inflates the pairwise baseline and consequently inflates raw M. The null 
z-scoring approach of Liardi et al.~\cite{liarinull} is the appropriate tool 
for isolating genuine higher-order collective structure from this pairwise 
confound, but the current implementation requires revision before the results 
can be reported with confidence. Finalised z-score results will be presented 
in the full dissertation.

\subsection{Representational Geometry}

PCA of the hidden state matrix confirms the geometric prediction from 
Section~\ref{sec:geometry}. Local-Hom networks have steeper PCA variance 
decay: more variance is concentrated in the first one or two principal 
components, indicating that dynamics across all neurons are strongly 
correlated and occupy a low-dimensional subspace. FittedHet networks show 
more gradual decay, requiring more components to explain the same fraction 
of total variance.

The Local-Hom network on the 4-class task concentrates 86.3\% of variance in
PC1 and reaches 90\% explained variance in two components; FittedHet-LogNorm
requires three. On the all-class task FittedHet networks require more components,
consistent with richer representational geometry under higher task complexity.

\begin{figure}[htb]
    \centering
    \includegraphics[width=\textwidth]{fig_pca_variance}
    \caption{PCA variance decay for all nine networks. Left: cumulative 
    explained variance ratio (EVR) against component index. Right: 
    per-component EVR on a log scale. Local-Hom networks (which share a 
    single timescale) show steeper decay with variance concentrated in 
    fewer components. FittedHet networks show more distributed variance 
    decay, consistent with per-neuron timescale diversity producing 
    higher-dimensional representations.}
    \label{fig:pca}
\end{figure}

When PCA dimensionality ($n_{90}$) is compared across all nine networks,
networks with higher representational dimensionality tend to show higher
collective synergy above the pairwise baseline.

\subsection{Membrane Potential versus Spike Trains}

Computing M-information from binary spike outputs as well as from membrane 
potentials is a planned component of the analysis, motivated by the observation 
that the thresholding nonlinearity may concentrate rather than destroy collective 
coding structure. However, this comparison depends on the correct extraction of 
binary spike tensors from the recurrent layer, which is part of the pipeline 
correction currently in progress (see Section~\ref{sec:zscore_plan}). 
A revised membrane-versus-spike comparison will be reported in the final 
dissertation once the corrected pipeline is validated.

\subsection{Summary of Key Preliminary Findings}

The analysis points to four inter-related findings:

\begin{enumerate}

    \item \textbf{The heterogeneity advantage varies with task complexity.} 
    The accuracy gap between FittedHet and Local-Hom architectures is small 
    on the binary task (+5pp) and large on the 4-class task (+10pp), 
    consistent with the hypothesis that parallel timescale channels are most 
    valuable when the task demands simultaneous integration across multiple 
    temporal dimensions.

    \item \textbf{Null z-scoring is essential for architecture comparison, but the pipeline requires correction before results can be reported.}
    The corrected implementation is described in Section~\ref{sec:zscore_plan}.

    \item \textbf{FittedHet networks produce higher-dimensional 
    representations.} PCA confirms that per-neuron timescale diversity leads 
    to a more distributed use of the representational space across all task 
    conditions, consistent with the theoretical predictions of 
    Wakhloo et al.~\cite{wakhloo2026} and Sussillo and Barak~\cite{sussillo2013}.

    \item \textbf{Spike trains carry comparable or higher collective
    information than membrane potentials.} Both representations are included
    in the corrected analysis pending spike-extraction validation.

\end{enumerate}


%%%%%%%%%%%%%%%%%%%%%%%%%%%%%%%%%%%%
\chapter{Project Plan}
\label{chap:plan}
%%%%%%%%%%%%%%%%%%%%%%%%%%%%%%%%%%%%

\section{Remaining Work}

The analysis described in Chapter~\ref{chap:achiev} establishes the core 
comparative results across nine network instances. The following areas represent 
the primary directions for the remainder of the project.

\paragraph{Null z-scoring pipeline: correction and re-validation.}
\label{sec:zscore_plan}
The most immediate priority is correcting and re-validating the null z-scoring
pipeline. A diagnostic audit identified that the original spike tensor extraction
incorrectly indexed the recurrent layer record tuple: index~0 (synaptic current) was
used instead of index~2 (binary spike output), invalidating all previously cached
spike z-score results. The membrane-potential results are unaffected. A second issue:
the original null model fit a Gaussian to continuous activations; for binary spike
data, a discrete null preserving Bernoulli spike statistics is theoretically
required~\cite{liarinull}. The revised pipeline replaces the continuous null generator
with a discrete model operating on binary spike transition probabilities, with the
temporal delay re-calibrated to a biologically motivated 15~ms window. Once
validated, the full subset sweep ($k \in \{2, 4, 8, 16, 32\}$) will be re-run for
all nine networks. The qualitative hypothesis that heterogeneous architectures show
higher genuine collective synergy above their pairwise
baseline~\cite{liarinull, liardiscalable} remains the central claim to be tested.

\paragraph{Extended analysis of the 2-class binary SHD task.}
The 2-class binary task warrants deeper investigation as a tractable reference
point: the maximum mutual information between a two-class stimulus and the network
output is exactly 1 bit, establishing a well-defined upper bound against which
M-information can be anchored. This makes it possible to assess how
\emph{efficiently} each architecture converts input temporal structure into
collective representations. Preliminary results reveal that architecture rankings
shift as the subset size $k$ grows, suggesting that binary discrimination is
achievable through multiple distinct coding strategies. The binary task is also the
most tractable setting for multi-seed replication, and 5--10 replicates per
architecture are planned to confirm that observed information signatures reflect
structural properties of the architecture family rather than random-seed sensitivity.

\paragraph{New benchmarking and visualisation of representational geometry.}
The current PCA analysis provides a global linear summary of representational
dimensionality. A richer geometric analysis will be applied to the 2-class and
4-class tasks: UMAP and t-SNE~\cite{mcinnes2018, vandermaaten2008} for nonlinear
class manifold visualisation; RSA~\cite{kriegeskorte2008} and
CKA~\cite{kornblith2019} to quantify cross-architecture representational similarity;
and the manifold capacity framework of Chung et al.~\cite{chung2018} for a
task-informed geometric capacity measure. Together these tools will establish whether
architectures with similar accuracy scores share a common representational structure,
and whether the FittedHet networks converge to the geometry of Repo-Learned-Het.

\paragraph{Multi-seed robustness.}
All current results are based on single trained instances of each 
architecture. To confirm that observed differences reflect structural 
properties of the architecture families rather than sensitivity to a single 
random seed, 5--10 additional replicates of each network type will be trained 
on the 2-class task. M z-scores and PCA dimensionality statistics will be 
reported as mean $\pm$ standard error across replicates, and the key 
cross-architecture comparisons will be assessed for statistical consistency.

\paragraph{Report preparation.}
The analysis currently exists in notebook and JSON form. Consolidating 
results into publication-quality figures, interpreting findings in the 
context of the literature, and producing a coherent final report narrative 
represents a substantial remaining task. Several figures (convergence 
diagnostics, membrane-vs-spike comparisons, global percentile curves) 
need final polish for the written report.

\section{Project Timeline}

Table~\ref{tab:timeline} summarises the planned schedule for the 
remaining project phases.

\begin{table}[htb]
    \centering
    \small
    \begin{tabular}{p{5.5cm}|p{3cm}|p{5.5cm}}
        \textbf{Phase} & \textbf{Target date} & \textbf{Deliverable} \\ \hline
        Background report submission
            & June 2026
            & This document \\ \hline
        Null z-scoring pipeline correction
            & June--July 2026
            & Corrected spike extraction, discrete null model, and
              validated z-score sweep across all nine networks \\ \hline
        Multi-seed robustness (2-class)
            & July 2026
            & Replicated z-score and PCA statistics with error bars \\ \hline
        2-class binary deep analysis
            & July 2026
            & Extended M z-score profiles, per-neuron contributions, and
              geometric analysis for all 2-class architecture pairs \\ \hline
        Architecture expansion training
            & July--August 2026
            & 20 new model checkpoints (10 homogeneous, 10 heterogeneous)
              trained on 2-class SHD task \\ \hline
        Manifold visualisation and CKA/RSA
            & August 2026
            & UMAP projections, RSA dissimilarity matrices, and CKA
              similarity scores across all architectures \\ \hline
        Figures and draft writeup
            & August--September 2026
            & Complete final report draft \\ \hline
        Final report submission
            & September 2026
            & Submitted thesis \\
    \end{tabular}
    \caption{Planned timeline for the remaining project phases.}
    \label{tab:timeline}
\end{table}

\section{Defining Success}

The project will be considered successful if it can provide clear, 
evidence-based answers to the following questions:

\begin{enumerate}

    \item Does the \emph{type} of temporal heterogeneity (fitted
    versus learned, log-normal versus log-uniform) produce measurably 
    different collective coding signatures in trained SNNs, as captured 
    by M z-scores and PCA dimensionality?

    \item Is the relationship between heterogeneity and collective synergy 
    task-dependent in the predicted direction: smallest advantage on the 
    binary task and largest on the most complex multi-class task?

    \item Can the geometric link between representational dimensionality 
    and synergistic collective coding be formally characterised, and does 
    the participation ratio~\cite{wakhloo2026} provide a useful quantitative 
    bridge between PCA eigenspectra and M-information?

    \item Are the results robust to training randomness, as verified 
    by multi-seed replication?

\end{enumerate}

\subsection*{Primary success criterion}

The primary success criterion is as follows: across the expanded suite of
approximately twenty networks trained on the 2-class SHD task, a statistically
significant majority of heterogeneous models should show higher M z-scores than their
homogeneous counterparts at matched population sizes and under the same
null-normalisation procedure~\cite{liarinull}. A result will be considered robust if
the M z-score advantage of heterogeneous models over Local-Hom holds at $p < 0.05$ on
the null test in at least seven of the ten heterogeneous architectures tested, and
replicates consistently across seeds and subset sizes. A further aim is demonstrating
that the null z-scoring and M-information pipeline is stable and interpretable across
diverse SNN architectures and task conditions, which would itself constitute a
methodological contribution given that this combination has not previously been applied
to trained SNN classification models.


%%%%%%%%%%%%%%%%%%%%%%%%%%%%%%%%%%%%
% Bibliography
%%%%%%%%%%%%%%%%%%%%%%%%%%%%%%%%%%%%
\bibliographystyle{unsrt}
\bibliography{references}

\end{document}
